\section{Расчёт выведения учебной двухступенчатой ракеты}
\label{sec:numerical-ascent}
% Generated by ascent_numerics.py; units documented in source.
\newcommand{\AscentTOne}{134{,}491}
\newcommand{\AscentRaise}{147{,}286}
\newcommand{\AscentCoast}{301{,}229}
\newcommand{\AscentCircularize}{83{,}044}
\newcommand{\AscentFinalTime}{666{,}050}
\newcommand{\AscentFinalMass}{12{,}245}
\newcommand{\AscentReserve}{3{,}245}
\newcommand{\AscentCharSpeed}{10{,}048}
\newcommand{\AscentGravityLoss}{2303{,}188}
\newcommand{\AscentDragLoss}{36{,}440}
\newcommand{\AscentSteeringLoss}{177{,}407}
\newcommand{\AscentMaxQ}{31{,}508}
\newcommand{\AscentMaxQAlpha}{0{,}723}
\newcommand{\AscentMaxLoad}{6{,}662}
\newcommand{\AscentPositionCheck}{1.20\cdot10^{-7}}
\newcommand{\AscentVelocityCheck}{3.32\cdot10^{-10}}
\newcommand{\AscentEccentricity}{8.28\cdot10^{-14}}
\newcommand{\AscentResidual}{4.66\cdot10^{-14}}
\newcommand{\AscentBalanceError}{-6.22\cdot10^{-12}}


\subsection{Исходные данные и программа}

Выведем нагрузку массой $5$~т на круговую орбиту высотой $650$~км.
Цель примера --- проследить, как программа тяги формирует конечную
орбиту. Это учебная ракета; расчёт выполняется в плоскости вокруг
невращающейся сферической Земли. Примем
\[
 \mu=398600{,}4418\ \text{км}^3/\text{с}^2,\quad
 R_E=6378{,}137\ \text{км},\quad g_0=9{,}80665\ \text{м}/\text{с}^2.
\]
На старте $\vect R_0=(R_E,0)^\top$, $\vect V_0=0$;
целевые значения $r_*=7028{,}137$~км и $V_*=7{,}5309329$~км/с.

\begin{table}[htbp]
\centering
\caption{Характеристики учебной ракеты.}
\label{tab:ascent-vehicle}
\renewcommand{\arraystretch}{1.12}
\begin{tabularx}{\textwidth}{@{}Xrr@{}}
\toprule
Параметр & Первая ступень & Верхняя ступень\\\midrule
Масса топлива, т & $160$ & $45$\\
Сухая масса, т & $16$ & $4$\\
Тяга, МН & $3{,}5$ & $0{,}8$\\
Удельный импульс, с & $300$ & $450$\\
\bottomrule
\end{tabularx}
\end{table}

Начальная масса равна $230$~т. Первая ступень расходует всё топливо
за $\AscentTOne$~с: к разделению остаётся $70$~т, после сброса
её сухой массы --- $54$~т. Верхний двигатель допускает выключение
и повторный запуск. Тяга и удельный импульс каждой ступени постоянны;
переходные процессы двигателя и поворота корпуса не моделируются.

Атмосферу зададим экспонентой
\begin{equation}
 \rho(H)=1{,}225\exp\left(-\frac{H}{7{,}2\ \text{км}}\right)
 \ \text{кг}/\text{м}^3,\qquad C_D=0{,}3,\quad S=10\ \text{м}^2.
 \label{eq:numerical-atmosphere}
\end{equation}
В сопротивлении скорость переводится в м/с.
Атмосфера неподвижна, поэтому воздушная и инерциальная скорости
совпадают; подъёмная сила не учитывается.

Программа состоит из шести участков:
\begin{enumerate}
\item Вертикальная тяга в течение $10$~с.
\item Линейное уменьшение $\theta$ от $90^\circ$ до $70^\circ$
за следующие $25$~с --- начальный разворот.
\item Тяга вдоль скорости, $\theta=\gamma$, до выгорания первой ступени.
\item После разделения --- первый активный участок верхней ступени
при $\theta=45^\circ$, формирующий переходную дугу.
\item Пассивный подъём при выключенном двигателе.
\item Трансверсальное довыведение, $\theta=0$, до получения
круговой скорости при $V_r=0$.
\end{enumerate}
Изменения направления на границах идеализируются мгновенными;
положение и скорость остаются непрерывными.

\courseplot[.98]{ascent-trajectory}{Траектория с выделением участков.
Дальность равна $R_E\varphi$, где $\varphi$ --- угол от радиуса точки
старта. Масштабы высоты и дальности различны.}

\subsection{Пристрелка длительностей}

Изменяем длительности первого верхнего активного участка,
пассивной паузы и довыведения:
\[
 \vect p=(\tau_{\mathrm a},\tau_{\mathrm c},\tau_{\mathrm d})^\top.
\]
Доступное топливо ограничивает сумму активных участков:
\[
 \tau_{\mathrm a}+\tau_{\mathrm d}\le248{,}231\ \text{с}.
\]
Используем невязку~\eqref{eq:residual} с допуском
$\norm{\vect{\mathcal R}}_2<10^{-11}$.
От начального приближения
$\vect p^{(0)}=(150,300,80)^\top$~с выполнены три ньютоновские
поправки с производными по центральным разностям.
Каждая оценка невязки включает полный прогноз от старта.
Получено
\begin{equation}
 \boxed{
 \tau_{\mathrm a}=\AscentRaise\ \text{с},\qquad
 \tau_{\mathrm c}=\AscentCoast\ \text{с},\qquad
 \tau_{\mathrm d}=\AscentCircularize\ \text{с}.}
 \label{eq:numerical-solution}
\end{equation}
Конечный момент $t_f=\AscentFinalTime$~с.
Здесь параметры округлены; для расчёта используются
их полные значения из сохранённых данных.

\courseplot[.88]{ascent-convergence}{Уменьшение нормы безразмерной
невязки за три поправки. Номер $0$ соответствует исходной программе.}

\subsection{Формирование орбиты}

После первого активного участка верхней ступени высота составляет
около $283$~км, а радиальная скорость --- $2{,}10$~км/с.
Далее ракета поднимается без тяги за счёт уменьшения
кинетической энергии.

\courseplot[.98]{ascent-state-history}{Высота, компоненты скорости
и углы на общей временной шкале. На пассивном участке линия $\theta$
не показана. При $V=0$ угол $\gamma$ не определён; начальное значение
$90^\circ$ обозначает вертикальный предел.}

Довыведение начинается до апогея: $H=640{,}845$~км,
$V_r=0{,}2973$~км/с и $V_t=3{,}9991$~км/с.
При трансверсальной тяге и пренебрежимо малом сопротивлении
\begin{equation}
 \dot V_r=\frac{V_t^2}{r}-\frac{\mu}{r^2}<0
 \quad\text{при}\quad V_t<\sqrt{\frac{\mu}{r}}.
 \label{eq:numerical-radial-deceleration}
\end{equation}
Поэтому положительная радиальная скорость уменьшается без
направленной вниз тяги. Пауза подобрана так, чтобы $V_r$
достигла нуля одновременно с получением круговой скорости.

\courseplot[.94]{ascent-final-insertion}{Довыведение: радиальная
скорость и дефицит поперечной скорости одновременно стремятся к нулю.}

Перед последним манёвром апогей переходной орбиты равен
$648{,}455$~км, а перигей лежит внутри Земли.
Отрицательная высота перигея описывает продолжение оскулирующей
коники; оставлять аппарат на такой траектории нельзя.
На пассивной дуге апсиды практически постоянны.
Довыведение поднимает и апогей, и перигей до $650$~км.
Таким образом, последний манёвр не сводится к импульсу точно
в апогее целевой высоты.

\courseplot[.98]{ascent-osculating-orbit}{Изменение апсид и удельной
механической энергии. Постоянный участок соответствует пассивной дуге;
в конце апогей и перигей совпадают.}

\FloatBarrier
\subsection{Топливо, нагрузки и потери}

Конечная масса равна $\AscentFinalMass$~т, включая $9$~т сухой
верхней ступени и нагрузки. Остаток топлива --- $\AscentReserve$~т.
На активных участках масса убывает непрерывно; разделение создаёт
её скачок. Уменьшение массы при постоянной тяге увеличивает $n_T$.

\courseplot[.98]{ascent-propulsion}{Масса, тяга и тяговая перегрузка.
Сброс массы при разделении не создаёт скачка скорости.}

После пристрелки проверены иллюстративные пределы
$q_{\mathrm{dyn}}\le50$~кПа и $n_T\le8$.
Максимумы на сетке с интервалом не более $1$~с составляют
$\AscentMaxQ$~кПа и $\AscentMaxLoad$ соответственно.
При тяге по продольной оси в начальном развороте
$|\alpha|_{\max}\simeq2{,}68^\circ$,
а $(q_{\mathrm{dyn}}|\alpha|)_{\max}=\AscentMaxQAlpha$~кПа$\cdot$рад.
На верхнем участке рассогласование тяги и скорости достигает
$20{,}4^\circ$, но напор уже не превышает $0{,}323$~кПа.
Это объясняет, почему одинаковое отклонение направления
создаёт разные аэродинамические нагрузки на разных высотах.
Отдельный расчёт прочности здесь не выполняется.

\courseplot[.98]{ascent-aeroloads}{Напор и произведение
$q_{\mathrm{dyn}}|\alpha|$ на первой ступени; угол выражен в радианах.
После перехода к тяге вдоль скорости $\alpha=0$.}

Скоростной баланс даёт
\begin{equation}
 \underbrace{10{,}047968}_{\Delta V_{\mathrm{th}}}
 -\underbrace{2{,}303188}_{\Delta V_g}
 -\underbrace{0{,}036440}_{\Delta V_D}
 -\underbrace{0{,}177407}_{\Delta V_{\mathrm{steer}}}
 =7{,}530933\ \text{км}/\text{с}.
 \label{eq:numerical-speed-budget}
\end{equation}
Гравитационный интеграл растёт и на пассивной восходящей дуге:
характеристическая скорость постоянна, а текущая скорость уменьшается.
Сопротивление в основном влияет на первую ступень;
рулевые потери накапливаются прежде всего при $\theta=45^\circ$.

\courseplot[.98]{ascent-losses}{Накопленные потери и различие
характеристической и текущей скоростей. Подъём потенциальной энергии
уже учтён в балансе через гравитационный интеграл.}

\FloatBarrier
\subsection{Влияние начального разворота}

Сохраним найденные длительности и изменим только начальный разворот:
$16^\circ$ или $24^\circ$ вместо $20^\circ$.
Расход топлива и характеристическая скорость останутся одинаковыми.

Крутая программа заканчивается на высоте $793$~км,
но её перигей около $-126$~км: большая высота не обеспечила
орбитальный полёт. Пологая программа заканчивается на высоте
$501$~км с отрицательной $V_r$ и получает эллиптическую орбиту
с апсидами $447$ и $1326$~км. Обе программы не выполняют
круговую цель.

\courseplot[.98]{ascent-program-comparison}{Разные результаты трёх
начальных разворотов при одинаковых последующих длительностях
и расходе топлива. Соседние программы повторно не пристреливались.}

\begin{explanation}{Результат примера}
Пристрелка обеспечила круговую орбиту $650$~км с положительным
остатком топлива и выполнением выбранных проверок напора и $n_T$.
Она не минимизировала расход и не максимизировала полезную нагрузку.
Сравнение программ показывает роль управления при одинаковом
запасе характеристической скорости.
\end{explanation}

Результат подтверждён повторным интегрированием с более строгими
допусками и проверками массового и скоростного балансов.
Расчёт и рисунки воспроизводятся скриптом
\texttt{10/images/ascent\_numerics.py}; параметры, траектории и
проверки сохранены рядом в JSON, CSV и NPZ.
